Source : Insee, comptes nationaux annuels base 2014 (séries arrêtées) jusqu'en 1995 ; Eurostat et Insee base 2020 ensuite. Calculs de l'auteur (dette/dynamique\_longue.py).\par Lecture : les prestations sociales (en espèces, et en nature de produits marchands comme les remboursements de soins) passent de 11,3 \% du PIB en 1960 à 17,9 \% en 1978 et 25,8 \% en 2025. Les administrations ont dégagé des excédents primaires presque chaque année jusqu'au milieu des années 1970, puis en 1979-1980, 1989-1991 et 1998-2001 : la hausse des prestations a longtemps été financée par celle des prélèvements. Le trait vertical marque le passage de la base 2014 à la base 2020.
